%% Included by 05-experiments.tex only when sections/misspec-numbers.tex exists (E08 done).
\begin{figure*}[t]
  \centering
  \includegraphics[width=\textwidth]{fig_misspec.pdf}
  \caption{\textbf{Where the structure stops helping.} Ratio of a radius-2 window
  head's error to EM--BP's against the misspecification strength, for four Gaussian
  departures from the chain with exact reference scores (rank-one global latent,
  long-range precision coupling, an AR(2) second coefficient at fixed lag-one correlation,
  site-dependent $\corr$), at two sample sizes (top row). Above the dotted line the misspecified
  chain estimator still wins. Bottom row: EM--BP's own error at both sizes (points,
  paired standard errors) against the closed-form bias $b(\varepsilon)$ of the
  Kullback--Leibler projection (dotted) and the finite-sample prediction
  $\sqrt{b(\varepsilon)^2+v/\nseq}$ (lines) with $v$ read from the well-specified cells
  (\cref{app:audit}).}
  \label{fig:misspec}
\end{figure*}
Every departure is Gaussian so that its score is exact, and every arm is scored against
that score; the network arm is a radius-2 window head under the recipe of
\cref{app:protocol}. The ordering degrades smoothly in each mechanism and the strength
at which it breaks even is a property of the mechanism, not a single number. A rank-one
global latent never reverses the ordering in the swept range: the ratio falls only from
$\misRatioGlobalEzeropzeroD$ to $\misRatioGlobalEonepzeroD$ across the swept range, up
to $\beta=\misEpsMaxGlobal$. An AR(2) second coefficient reverses it between
$\misBreakLoArtwoD$ and $\misBreakHiArtwoD$ at both sizes, and a site-dependent $\corr$
between $\misBreakLoSitevarD$ and $\misBreakHiSitevarD$ at $\nseq=2048$. Long-range
precision coupling does not reverse it anywhere we swept: the ratio decays to
$\misRatioLongrangeEzeropfourD$ at $\gamma=\misEpsMaxLongrange$, a tie rather than a
crossing. We report brackets and not fitted crossings, because with $\misSeeds$ seeds
the paired log-ratio is resolved to a few thousandths and the crossing sits between grid
points either way.

The two sizes separate the statistical gain of the structure at small $\nseq$ from its
asymptotic bias at large $\nseq$, and the bias has a closed form. Since every truth is
Gaussian and the fitted innovation is Gaussian to three decimals, the $\nseq\to\infty$
limit of EM--BP is the Kullback--Leibler projection of the noised truth onto the
Gaussian AR(1) family under the one-view protocol, and its relative score error
$b(\varepsilon)$ is a trace formula (\cref{app:audit}). It reproduces the measured
EM--BP error at $\nseq=2048$ to within $\audBiasMaxAbsDevLagoneD$ across every
mechanism and strength, and $\sqrt{b(\varepsilon)^2+v/\nseq}$ with $v$ read from the
well-specified cells predicts both sizes: at zero strength the EM--BP error falls from
$\misArtwoEmbpEzeropzeroC$ to $\misArtwoEmbpEzeropzeroD$ between $\nseq=512$ and $2048$,
which is estimation; at the largest AR(2) strength it moves from
$\misArtwoEmbpEzeropfourC$ to $\misArtwoEmbpEzeropfourD$, which is bias. The order of the
bias in the strength explains which mechanisms bite: $b\propto\varepsilon$ for the
AR(2) coefficient, the site-dependent $\corr$ and the long-range coupling (local
exponents $\audBiasOrderArtwo$, $\audBiasOrderSitevar$, $\audBiasOrderLongrange$), which
perturb of order $\nsites$ directions of the covariance, and $b\propto\beta^2$ for the
rank-one latent ($\audBiasOrderGlobal$), which is even in the latent and perturbs one
direction of $\nsites$. Past the break-even the chain estimator has converged to a wrong
law and more data no longer buys anything back---exactly the behaviour \cref{sec:theory}
attributes to a model whose bias, not whose variance, is binding.
